\documentclass{article}

% Language setting
\usepackage[english]{babel}

% Set page size and margins
\usepackage[letterpaper,top=2cm,bottom=2cm,left=3cm,right=3cm,marginparwidth=1.75cm]{geometry}

\usepackage{amsmath}
\usepackage{graphicx}
\usepackage[colorlinks=true, allcolors=blue]{hyperref}

\title{Reference Manual for TwoDSensorModule}
\author{Hojae Lee}

\begin{document}
\maketitle

\section{Introduction}

This document is a reference manual to \href{https://github.com/thousandbrainsproject/feat.2d_sensor_module}{\texttt{TwoDSensorModule}}. The document will cover results and explanation of how and why certain algorithms are implemented, and also point out limitations and areas for improvement. The manual will be detailed enough for someone who is interested in improving any part of the TwoDSensorModule to have a good understanding on why certain things were chosen and get a guide for where to go.

The implementation of TwoDSensorModule was roughly divided into two phases:

\begin{enumerate}
  \item Detecting dominant edge in a patch
  \item Implementing 2D Movement
\end{enumerate}

\section{Part I: Computing Orientation of Dominant Edge in Patch}

Key functions:
\begin{enumerate}
  \item \texttt{compute\_edge\_features}
  \item \texttt{edge\_angle\_to\_2d\_pose}
  \item \texttt{is\_geometric\_edge}, \texttt{gradient\_to\_tangent\_angle}
\end{enumerate}

The Problem. This feature was initially posed as a edge detection problem, however, in reality a much more complex problem. While functions that extract surface normal rely on extracting the value at the center of the 64 x64 patch, we are not interested in whether an edge exists at the center, but rather whether there is a dominant edge within the patch and what its orientation is. What constitutes as "dominant" and how much it should be "center-focused"

The same image at different zoom levels can have different edge orientations.

Below figure demonstrates some examples and desired properties of this problem.

% Add figure here

The Solution. The main algorithm that performs dominant edge detection as well as computing its orientation, edge strength, and coherence is implemented in \texttt{compute\_edge\_features}

Note that we have explored well-known edge detection methods such as Canny Edge Detector (cv2.CannyEdge) and

Additional Details.

Note that by itself, \texttt{compute\_edge\_features}
Helper Functions and brief explanations
\begin{enumerate}
  \item \texttt{GaussianBlurRGB} in \texttt{transforms.py}
\end{enumerate}

Part I of implementing TwoDSensorModule

\section{Part II: Implementing 2D Movement}

Key Functions:
\begin{enumerate}
  \item \texttt{TangentFrame}
  \item \texttt{arc\_from\_projection}
\end{enumerate}

Helper Functions:
\begin{itemize}
  \item \texttt{directional\_curvature}
  \item \texttt{arc\_length\_corrected\_displacement}
\end{itemize}

\subsection{\texttt{TangentFrame}}

\textbf{The Problem.}
As the 2D sensor moves across a curved 3D surface, we need to express
sensor displacement in a consistent 2D coordinate frame. On a flat
surface a fixed $(u, v)$ pair suffices, but on a curved surface the
tangent plane rotates with the curvature. Naively accumulating
displacement vectors in a fixed world frame causes the tracked position
to drift off the surface.

\textbf{The Solution.}
\texttt{TangentFrame} maintains a right-handed orthonormal basis
$(u, v, n)$, where $n$ is the current surface normal, $u$ is the
horizontal tangent direction, and $v$ is the vertical tangent direction.
As the sensor moves to a new surface point, the method
\texttt{transport(new\_normal)} updates the basis via
\emph{parallel transport}---the unique minimum-rotation update that
keeps $u$ and $v$ in the new tangent plane without introducing
extraneous twist. This is the discrete analogue of parallel-transporting
a vector along a geodesic on a Riemannian manifold
(\texttt{https://en.wikipedia.org/wiki/Parallel\_transport}).

\textbf{Initialization.}
Given an initial surface normal $n$, the constructor builds $u$ and $v$
as follows. It picks an arbitrary reference axis $a = [0, 1, 0]$,
falling back to $a = [0, 0, 1]$ if $n$ is nearly parallel to $a$
($|\cos\theta| > 1 - \epsilon$). Then:
\begin{align*}
  u &= \mathrm{normalize}(a \times n) \\
  v &= \mathrm{normalize}(n \times u)
\end{align*}
The cross products ensure $u$ lies in the tangent plane and $(u, v, n)$
forms a right-handed frame.

\begin{figure}[h]
  \centering
  \includegraphics[width=0.8\textwidth]{../imgs/tangent_frame_init.png}
  \caption{Initialization of the orthonormal tangent frame $(u, v, n)$ from a surface normal $n$.}
\end{figure}

\textbf{Parallel Transport.}
When \texttt{transport(new\_normal)} is called with a new surface normal
$n'$, the update proceeds in three steps:
\begin{enumerate}
  \item Compute the rotation $R$ that maps $n \to n'$ via the
    axis-angle representation:
    \[
      \text{axis} = \mathrm{normalize}(n \times n'), \quad
      \theta = \arccos(\mathrm{clip}(n \cdot n', -1, 1))
    \]
  \item Apply $R$ to $u$: this keeps $u$ as close as possible to its
    previous direction while moving it into the new tangent plane.
  \item Re-orthonormalize to suppress floating-point drift accumulated
    over many steps:
    \[
      u \leftarrow \mathrm{normalize}(u - (u \cdot n') \, n'),
      \quad v \leftarrow n' \times u
    \]
\end{enumerate}
Step 3 is necessary because repeated small rotations accumulate
numerical error; projecting $u$ back onto the tangent plane and
recomputing $v$ from scratch keeps the frame exactly orthonormal.

\begin{figure}[h]
  \centering
  \includegraphics[width=0.6\textwidth]{../imgs/tangent_frame_transport.png}
  \caption{Parallel transport of the tangent frame from normal $n$ to a new normal $n'$.}
\end{figure}

\subsection{\texttt{arc\_from\_projection} and Arc-Length Correction}

\textbf{The Problem.}
When the sensor moves along a curved surface, the displacement we observe is a
chord in the tangent plane, not the true distance traveled along the curve.
On a convex surface the sensor must travel further than the straight-line
projection suggests; without correction, accumulated 2D positions are
systematically compressed relative to the actual surface path.

\textbf{The Solution.}
At each surface point we fit an \emph{osculating circle}: the unique circle
that is tangent to the surface and has the same normal curvature $k$ in the
direction of movement. The osculating circle is the best local circular
approximation to the curve, in the sense that it agrees with the curve up to
second order. The approximation holds well when curvature is approximately
constant over the displacement, but becomes inaccurate for surfaces with
rapidly varying curvature (e.g., a sharp ridge crossed in a single large step).

Under this assumption the relationship between tangent-plane chord $p$ and
arc length $s$ is:
\[
  p = \frac{\sin(k \cdot s)}{k}
  \quad\Longrightarrow\quad
  s = \frac{\arcsin(k \cdot p)}{k}
\]
This is implemented in \texttt{arc\_from\_projection(tangent\_projection, curvature)}.
Two guard conditions keep the formula numerically safe:
\begin{itemize}
  \item $|k \cdot p| < 0.001$: the surface is nearly flat over the step, so
        $\arcsin(k p) \approx k p$ and no correction is applied.
  \item $|k \cdot p| \geq 1.0$: the chord exceeds the radius of curvature,
        meaning the sensor has moved past the most extreme visible point under
        the constant-curvature assumption. The correction is skipped and a
        warning is logged; the proper remedy is a policy that takes smaller 3D
        steps.
\end{itemize}
The formula is sign-preserving (\texttt{np.copysign}) and works for both
convex ($k > 0$) and concave ($k < 0$) surfaces because the arc-to-chord
geometry on a circle is the same regardless of the sign of curvature.

\textbf{Helper: \texttt{directional\_curvature}.}
Before correcting displacement we need the normal curvature $k$ in the
direction of movement. Given principal curvatures $k_1, k_2$ and their
directions $\hat{d}_1, \hat{d}_2$, Euler's curvature formula gives:
\[
  k(\theta) = k_1 \cos^2\theta + k_2 \sin^2\theta
\]
where $\theta$ is the angle between the movement direction and $\hat{d}_1$.
\texttt{directional\_curvature} validates that $\hat{d}_1$ and $\hat{d}_2$
are orthonormal and that the movement direction lies in their span before
applying this formula.

\textbf{Helper: \texttt{arc\_length\_corrected\_displacement}.}
This function ties the two helpers together. Given tangent-frame displacements
$(du, dv)$ along basis vectors $(u, v)$, it:
\begin{enumerate}
  \item Calls \texttt{directional\_curvature} once for $u$ and once for $v$
        to obtain $k_u$ and $k_v$ from the observed principal curvatures and
        their pose vectors.
  \item Passes each chord-displacement through \texttt{arc\_from\_projection}
        to obtain arc-length-corrected displacements $(\hat{u}, \hat{v})$.
\end{enumerate}
The result is a pair of corrected 2D displacements that accurately reflect the
distance the sensor traveled along the surface in each tangent direction.

\section{Putting it Together in TwoDSensorModule}

Currently, the \texttt{TwoDSensorModule} is written similarly to \texttt{CameraSM}, with two major differences:
\begin{enumerate}
  \item After extracting 3D morphological features, it extracts edge features
\end{enumerate}


\end{document}
